% Pursuing Stacks (cotes 134-2 to 134-8): archive cote 134-4, site page 24 (PDF page 25), typescript page 112.
% typescript. Draft compared: ps3.tex:1177-1245.
% AI-transcribed by Claude Opus 5.5 (claude-opus-5-5, adaptive thinking, effort high), Message Batches API, batch ps2 (msgbatch_01NVwr1Yuo8qJf5rykeGBXPK), 2026-11-12; not yet checked by eye.
\scanpage{134-4/24}
\mspage{112}

A) \uline{Homotopy relation between maps}. As a \uline{type of structure}, a
homotopy relation between maps in M is a subset
\[ R_h \subset F\!\not{l}(M)\mathrm{x}F\!\not{l}(M) \]
namely a relation in the set $F\!\not{l}(M)$ or arrows of M , the \uline{basic assumption}
being that whenever f and g are \hand{``}homotopic\hand{''} arrows, then they have same
source, and same target. Thus the data $R_h$ is equivalent to giving a
``homotopy relation'' in any one of the sets Hom(X,Y), with X and Y objects
in M. The relevant \uline{saturation condition} is twofold:

a) The relation $R_h$ is an equivalence relation, or equivalently, the
corresponding relations in the sets Hom(X,Y) are \unsure{equivalehce|equivalence} relations;

b) stability under composition: if f and g are homotopic, then so
are vf and vg , and so are fu and gu , for any arrow v or u such that
the relation makes sense.

When these conditions are satisfied, we'll say we got a \uline{homotopy}
\uline{relation} \xout{\unsure{in M}} between maps of M. This relation \inserted{between f and g} will be denoted by a sym-%
bol like
\[ f \underset{h}{\hand{\sim}} g \qquad . \]
If the basic assumption is satisfied, but not the saturation condition,
there is an evident way of ``saturating'' the given relation, getting
\deleted{a saturated} one $\overline{R}_h$ which is saturated i.e. a homotopy relation in M
(in fact, the smallest one containing $R_h$).

Given a homotopy relation $R_h$ , we denote by \xout{Hom}
\[ \mathrm{Hom}(X,Y)_h \]
the corresponding quotient sets of the sets Hom(X,Y), they compose in an
evident way, so as to give rise to a \xout{functor} category $M_h$ having same
objects as M, and to a canonical functor
\[ M \longrightarrow M_h \]
which is the identity on objects, and surjective on arrows. We may
view thus $M_h$ as a \uline{quotient category} of M, having same objects as M.
Clearly, $R_h \hand{\longmapsto} M_h$ is a bijective correspondance between the set of homo-%
topy relations in M, and the set of quotient categories of M satisfying
the aforesaid property. By abuse of language, we may even consider that
considering a homotopy relation in M, amounts to the same as giving a
functor $M \hand{\rightarrow} M_h$ \inserted{from M} which is bijective on objects and surjective on arrows.

When we got a homotopy relation $R_h$ in M, we deduce a notion of
\uline{homotopisms}
\[ W_h \subset F\!\not{l}(M) \ , \]
namely those arrows in M which \xout{xxx} become isomorphisms in $M_{\hand{h}}$ . Also we

\margin{24}
